\subsection*{Composing a description from pictured facts}
A reader trained on 33,525 human descriptions learned to write a sentence one glyph at a time. For the first Chinese test card with an editable family policy, its answer described a pub near the Rainbow Vegetarian Caf\'e, the cuisine, location and exclusion of families. Changing the single policy field yielded a new sentence welcoming families while preserving the other facts (Fig.~\ref{fig:semantic-page}a). All 1,847 selected Chinese-card sentences differed from the training strings. The first test card also exposed an invented price. On the last card, the writer named the wrong nearby establishment.

Greedy writing, likelihood beam search and optical readback recovered the exact literal fact set on 506, 576 and 657 Chinese cards, respectively. The detector accepted only 1,729 of 4,693 human references under this criterion. Price and rating patterns frequently identify two descriptions of the same scale. Equating the corresponding bands gave 1,258, 1,410 and 1,606 complete answers and 4,396 accepted human references. The matched byte reader selected 1,610 complete answers on Chinese cards under this convention. The readback writer got the direction of the changed fact right in both sentences on 1,286 of 1,291 pairs, while both entire answers passed the literal test on 418 pairs. Under cyclic exchange of the input records, zero of its 1,847 original Chinese-card answers passed literally.

Holding the eight completed sentences and the length normalization fixed isolated the readback term. Switching it on raised complete replies from 1,453 to 1,606 under the merged criterion, with 179 gains and 26 losses. At least one candidate passed on 1,668 records. Consequently 179 records lacked a passing candidate and on another 62 the selector missed one. In an unblinded, model-assisted reading of twenty evenly spaced records, the greedy and selected sentences conveyed all requested facts without adding another on 13 and 17, respectively. Wording remained awkward in several faithful answers. On a separate development set with Chinese and English rows interleaved within each card, the frozen reader passed 1,304 of 1,484 records with equivalent bands and 246 literally.

The optical writer and listener used 11.51 and 1.49 kJ respectively of on-device training energy. The matched byte writer used 11.72 kJ and shared the listener. These small differences do not establish a computational advantage. The writer learns English prose in a finite semantic domain. It has a fixed output alphabet and structured inputs. The illustrated dictionary preserves forms outside this alphabet by retrieving their recorded ink.
